\documentclass[10pt,a4paper]{amsart}
\usepackage[margin=19mm]{geometry}
\usepackage{amsmath,amssymb,mathtools}
\usepackage[scr=rsfs,cal=euler]{mathalfa}
\usepackage{xfrac}
\usepackage[unicode,hidelinks,bookmarks=false]{hyperref}
\newcommand{\ie}{i.e.\ }
\newcommand{\cf}{cf.\ }
\newcommand{\resp}{respectively\ }
\newcommand{\Subsection}{\subsection}
\providecommand{\nobreakdash}{\nobreak\hbox{-}}
\setlength{\emergencystretch}{2em}
\multlinegap0pt
\usepackage{amssymb}
\usepackage{mathtools}
\newtheorem{theorem}{Theorem}
\newtheorem{proposition}{Proposition}
\newcommand{\Z}{\mathbb Z}
\newcommand{\ee}{\varepsilon}
\title{A negative answer to the Erd\H{o}s--S\'ark\H{o}zy question}
\author{Simone Costa}
\date{Source: arXiv:2609.06303v1 (CC BY 4.0)}
\begin{document}
\maketitle

\noindent\emph{Source and licence.} A negative answer to the Erd\H{o}s--S\'ark\H{o}zy question. Version arXiv:2609.06303v1 (CC BY 4.0), distributed by arXiv under the Creative Commons Attribution 4.0 International licence, http://creativecommons.org/licenses/by/4.0/, as recorded in the arXiv licence metadata for this exact version and shown on the abstract page https://arxiv.org/abs/2609.06303v1. Full licence notice: \texttt{licenses/arxiv-2609.06303-CC-BY-4.0.txt}.\par\smallskip

\noindent\textit{Editorial scope.} Counted unit: the article's theorem thm:main (source0.tex lines 91--97). The article's complete original proof of it is delivered with it (source0.tex lines 163--244, one \texttt{proof} environment of 82 lines, which the article places away from the statement). No mathematics is written, repaired or re-derived: the statement and the proof are the article's own lines, in the article's own notation, and no retained line is substituted or re-flowed.  Retained uncounted as the source-local context the proof needs: lines 111--119; lines 123--145.  Nothing was omitted from any retained span, so the package carries no omission record.  No retained line contains a citation, so the wrapper carries no bibliography and no bibliography entry.  No retained line contains a figure environment, an image-inclusion command or drawing code, so no image is delivered and the package declares no asset dependency.  Editorial formatting only, in addition: an AMS \texttt{amsart} preamble that restates the article's own declarations and macros used by the retained lines, namely the environments \texttt{theorem}, \texttt{proposition} on separate counters, together with the standard packages those lines need.  The retained environments print in this document as Theorem 1, Proposition 1 in order of appearance, whereas the article prints the same items with the numbering of its own full document.  The complete native source of the article as distributed in the arXiv e-print for this exact version is bundled unchanged as \texttt{sources/209520/source0.tex}.
\begin{proposition}[Korsky]\label{prop:korsky}
For every $m\ge1$,
\[
 g_3(m)=\min\left\{\max_{1\le r\le m}b_r:
 (b_1,\ldots,b_m)\in\Z_{>0}^m,\ 
 \varepsilon\longmapsto\sum_{r=1}^m\varepsilon_r b_r
 \text{ is injective on }\{0,1,2\}^m\right\}.
\]
\end{proposition}

\begin{proposition}[Consequence of the OpenAI construction]\label{prop:openai}
For every $K>0$ there exist integers $n\ge1$, $R\ge1$, $D\ge1$, and $E\ge0$, together with integer coefficients
\[
 a_0(t),\ldots,a_n(t)\qquad(t\in\Z),
\]
such that:
\begin{enumerate}
\item $KD<R^n$;
\item $a_i(t)/t^n\to D$ as $t\to\infty$ for every $0\le i\le n$;
\item for all sufficiently large positive integers $t$, all $a_i(t)$ are positive;
\item if $t,Q$ are positive integers satisfying
\[
 Q\le (t-E)R,
\]
then the map
\[
 \{0,\ldots,Q-1\}^{n+1}\longrightarrow\Z,
 \qquad
 (x_0,\ldots,x_n)\longmapsto\sum_{i=0}^n a_i(t)x_i
\]
is injective.
\end{enumerate}
\end{proposition}

\begin{theorem}\label{thm:main}
For every $\ee>0$ there is an integer $d\ge2$ such that
\[
 g_3(d\ell)\le \ee\,3^{d\ell}
\]
for all sufficiently large integers $\ell$.
\end{theorem}

\begin{proof}[Proof of Theorem~\ref{thm:main}]
Fix $\ee>0$. Choose $K>1/(3\ee)$ and apply Proposition~\ref{prop:openai}. Keep the resulting integers $n,R,D,E$ fixed, and put
\[
 d=n+1.
\]
For a positive integer $\ell$, set
\[
 Q_\ell=3^\ell,
 \qquad
 t_\ell=E+\left\lceil\frac{3^\ell}{R}\right\rceil.
\]
Then
\[
 Q_\ell=3^\ell\le(t_\ell-E)R.
\]
Since $t_\ell\to\infty$, for all sufficiently large $\ell$ the integers
$a_0(t_\ell),\ldots,a_n(t_\ell)$ are positive, and Proposition~\ref{prop:openai}(4) gives that
\begin{equation}\label{eq:box3}
 (x_0,\ldots,x_n)\longmapsto
 \sum_{i=0}^n a_i(t_\ell)x_i
\end{equation}
is injective on $\{0,\ldots,3^\ell-1\}^{n+1}$.

Define
\[
 A_\ell=
 \left\{3^j a_i(t_\ell):0\le i\le n,\ 0\le j<\ell\right\}.
\]
We claim that the map
\[
 \{0,1,2\}^{(n+1)\ell}\longrightarrow\Z,
 \qquad
 (\delta_{i,j})\longmapsto
 \sum_{i=0}^n\sum_{j=0}^{\ell-1}\delta_{i,j}3^j a_i(t_\ell)
\]
is injective. Suppose that two choices of coefficients $\delta_{i,j},\delta'_{i,j}\in\{0,1,2\}$ give the same sum, and set
\[
 x_i=\sum_{j=0}^{\ell-1}\delta_{i,j}3^j,
 \qquad
 x'_i=\sum_{j=0}^{\ell-1}\delta'_{i,j}3^j.
\]
Then $0\le x_i,x'_i\le3^\ell-1$ and
\[
 \sum_{i=0}^n a_i(t_\ell)x_i
 =\sum_{i=0}^n a_i(t_\ell)x'_i.
\]
Injectivity in~\eqref{eq:box3} gives $x_i=x'_i$ for every $i$. Uniqueness of base-three expansion then gives $\delta_{i,j}=\delta'_{i,j}$ for every $i,j$. Thus the map is injective. Taking coefficient arrays with one entry equal to $1$ and all the others equal to $0$ shows in particular that the $(n+1)\ell$ integers $3^j a_i(t_\ell)$ are pairwise distinct, so
\[
 |A_\ell|=(n+1)\ell=d\ell.
\]
Proposition~\ref{prop:korsky} therefore gives
\begin{equation}\label{eq:g3bound}
 g_3(d\ell)\le \max A_\ell.
\end{equation}

It remains to estimate the largest element. Since the number of coefficients is fixed and
$a_i(t)/t^n\to D$ for every $i$,
\[
 \max_{0\le i\le n}a_i(t_\ell)=(D+o(1))t_\ell^n
 \qquad(\ell\to\infty).
\]
Moreover
\[
 t_\ell=\frac{3^\ell}{R}+O(1),
 \qquad
 t_\ell^n=\frac{3^{n\ell}}{R^n}(1+o(1)).
\]
Hence
\begin{align*}
 \max A_\ell
 &\le 3^{\ell-1}\max_i a_i(t_\ell)\\
 &=\left(\frac{D}{3R^n}+o(1)\right)3^{(n+1)\ell}.
\end{align*}
By Proposition~\ref{prop:openai}(1),
\[
 \frac{D}{3R^n}<\frac1{3K}<\ee.
\]
Combining this with~\eqref{eq:g3bound} yields, for all sufficiently large $\ell$,
\[
 g_3(d\ell)\le\ee\,3^{d\ell}.
\]
\end{proof}

\end{document}
